% !PRINTED 235
\weberpage{235}
%<<CONT
woraus sich für $\psi(n)$ der Ausdruck
\[
\psi(n) = n\,\Pi\Bigl(1+\frac{1}{p}\Bigr)
\]
ergibt, wie in §\,63\footnote{Vgl. \gesperrt{Dedekind}: Über die
elliptischen Modulfunktionen. Journal f. Mathematik, Bd.\ 83.}.

Nach (7) ist nun in den Formeln (6) für $n\omega'$ zu setzen:
\[
\frac{\displaystyle \gamma' + \delta'\frac{c+\partial\omega}{a}}{\displaystyle \alpha' + \beta'\frac{c+\partial\omega}{a}},
\]
und es lassen sich die linearen Transformationen der $f$-Funktionen
[§\,40, (4), (8), (11)] anwenden. Wir nehmen dabei

\weq{12}{c \equiv 0\ (\mathrm{mod}\ 16),}

so daß nach (3) und (8):
\[
\begin{pmatrix} \alpha', & \beta' \\ \gamma', & \delta' \end{pmatrix} \equiv \begin{pmatrix} a^3, & 0 \\ 0, & a \end{pmatrix}\ (\mathrm{mod}\ 16).
\]

Bezeichnen wir mit $\varrho$ die dritte Einheitswurzel

\weq{13}{\varrho = e^{-\frac{2\pi i}{3}[\alpha'(\gamma'-\beta')-(\alpha'^2-1)\beta'\delta']}\, e^{\frac{2n\pi i}{3}[\alpha(\gamma-\beta)-(\alpha^2-1)\beta\delta]},}

so erhalten wir
\weq{14}{\begin{aligned}
P_{00} &= \varrho\,\frac{f\left(\dfrac{c+\partial\omega}{a}\right)}{f(\omega)^n},\\
P_{01} &= \varrho\left(\frac{2}{a}\right)\frac{f_1\!\left(\dfrac{c+\partial\omega}{a}\right)}{f_1(\omega)^n},\\
P_{10} &= \varrho\left(\frac{2}{\partial}\right)\frac{f_2\!\left(\dfrac{c+\partial\omega}{a}\right)}{f_2(\omega)^n}.
\end{aligned}}

Der Quotient

\weq{15}{\frac{P_{10}}{P_{00}} = \left(\frac{2}{\partial}\right)\frac{f_2\left(\dfrac{c+\partial\omega}{a}\right)}{f\left(\dfrac{c+\partial\omega}{a}\right)}\left(\frac{f\omega}{f_2\omega}\right)^n}

gibt nach §\,54, (3), wenn man
\[
u(\omega) = \sqrt[4]{k}
\]
setzt, die Größen
%CONT>>
